\chapter{Кривые второго порядка, их геометрические свойства.}

\section{Общее уравнение кривой второго порядка}

\begin{defn}
\textit{Кривой второго порядка} называется множество точек плоскости, координаты которых удовлетворяют уравнению вида
\begin{equation}\label{21.0.general}
  Ax^2 + Bxy + Cy^2 + Dx + Ey + F = 0,
\end{equation}
где $A, B, C, D, E, F \in \bbR$ и хотя бы один из коэффициентов $A, B, C$ отличен от нуля.
\end{defn}

\begin{notion}
Уравнение~\eqref{21.0.general} при помощи поворота и параллельного переноса системы координат может быть приведено к одному из канонических видов: эллипса, гиперболы, параболы, а также вырожденных случаев (пара прямых, точка, пустое множество).
\end{notion}

\section{Эллипс}

\begin{defn}
\textit{Эллипсом} называется геометрическое место точек плоскости, сумма расстояний от каждой из которых до двух фиксированных точек $F_1$ и $F_2$, называемых \textit{фокусами}, есть величина постоянная $2a > 0$:
\begin{equation}\label{21.1.def}
  \rho(P, F_1) + \rho(P, F_2) = 2a.
\end{equation}
\end{defn}

Пусть фокусы расположены симметрично относительно начала координат: $F_1(-c, 0)$ и $F_2(c, 0)$, где $c \geq 0$. Обозначим $b^2 = a^2 - c^2 \geq 0$. Тогда из определения~\eqref{21.1.def} выводится

\begin{equation}\label{21.1.canon}
  \frac{x^2}{a^2} + \frac{y^2}{b^2} = 1, \quad a \geq b > 0
\end{equation}
--- \textit{каноническое уравнение эллипса}.

\begin{proof}
Пусть $r_1 = \rho(P,F_1)$, $r_2 = \rho(P,F_2)$. Из условия $r_1 + r_2 = 2a$ следует $r_1 = 2a - r_2$. Возводя в квадрат расстояния:
\[
r_1^2 = (x+c)^2 + y^2, \quad r_2^2 = (x-c)^2 + y^2.
\]
Вычитая: $r_1^2 - r_2^2 = 4cx$, откуда $r_1 - r_2 = \dfrac{2cx}{a}$. Совместно с $r_1 + r_2 = 2a$ получаем $r_1 = a + \dfrac{cx}{a}$, $r_2 = a - \dfrac{cx}{a}$. Подставляя в $r_2^2 = (x-c)^2+y^2$ и упрощая при $b^2 = a^2 - c^2$:
\[
\left(1-\frac{c^2}{a^2}\right)x^2 + y^2 = a^2 - c^2
\;\Longrightarrow\;
\frac{x^2}{a^2}+\frac{y^2}{b^2}=1. \qedhere
\]
\end{proof}

\begin{defn}
\textit{Эксцентриситетом} эллипса называется величина
\begin{equation}\label{21.1.exc}
  \varepsilon = \frac{c}{a} \in [0, 1).
\end{equation}
При $\varepsilon = 0$ ($c=0$, $a=b$) эллипс обращается в \textit{окружность} радиуса $a$.
\end{defn}

\begin{defn}
\textit{Директрисами} эллипса называются прямые
\begin{equation}\label{21.1.dir}
  x = \pm \frac{a}{\varepsilon} = \pm \frac{a^2}{c}.
\end{equation}
\end{defn}

\begin{thm}[Фокально-директрисное свойство эллипса]
Для любой точки $P$ эллипса отношение расстояния до фокуса к расстоянию до соответствующей директрисы постоянно и равно эксцентриситету:
\begin{equation}\label{21.1.focal}
  \frac{r_{1,2}}{d_{1,2}} = \varepsilon,
\end{equation}
где $d_{1,2}$ --- расстояние от точки до соответствующей директрисы.
\end{thm}

\begin{proof}
Для фокуса $F_2(c,0)$ и директрисы $x = \dfrac{a^2}{c}$ имеем $r_2 = a - \dfrac{cx}{a}$, $d_2 = \dfrac{a^2}{c} - x$. Тогда
\[
\frac{r_2}{d_2} = \frac{a - \tfrac{cx}{a}}{\tfrac{a^2}{c}-x} = \frac{\tfrac{a^2-cx}{a}}{\tfrac{a^2-cx}{c}\vphantom{\tfrac{A}{A}}} = \frac{c}{a} = \varepsilon. \qedhere
\]
\end{proof}

\section{Гипербола}

\begin{defn}
\textit{Гиперболой} называется геометрическое место точек плоскости, модуль разности расстояний от каждой из которых до двух фиксированных точек $F_1$ и $F_2$ (\textit{фокусов}) постоянен и равен $2a > 0$:
\begin{equation}\label{21.2.def}
  \abs{\rho(P,F_1) - \rho(P,F_2)} = 2a.
\end{equation}
\end{defn}

Расположим фокусы аналогично эллипсу: $F_{1,2}(\mp c, 0)$, $c > a > 0$. Введём $b^2 = c^2 - a^2 > 0$. Аналогично выводу канонического уравнения эллипса получим

\begin{equation}\label{21.2.canon}
  \frac{x^2}{a^2} - \frac{y^2}{b^2} = 1
\end{equation}
--- \textit{каноническое уравнение гиперболы}.

\begin{defn}
\textit{Эксцентриситетом} гиперболы называется величина
\begin{equation}\label{21.2.exc}
  \varepsilon = \frac{c}{a} > 1.
\end{equation}
\end{defn}

\begin{thm}[Асимптоты гиперболы]
Гипербола~\eqref{21.2.canon} имеет две асимптоты:
\begin{equation}\label{21.2.asymp}
  y = \pm \frac{b}{a} x,
\end{equation}
то есть расстояние от точки гиперболы до каждой из этих прямых стремится к нулю при $\abs{x} \to \infty$.
\end{thm}

\begin{proof}
Расстояние от точки $\left(x_0,\, \tfrac{b}{a}\sqrt{x_0^2-a^2}\right)$ гиперболы до прямой $bx - ay = 0$:
\[
\rho = \frac{\abs{bx_0 - a \cdot \tfrac{b}{a}\sqrt{x_0^2-a^2}}}{\sqrt{a^2+b^2}}
= \frac{b\abs{x_0 - \sqrt{x_0^2-a^2}}}{c}
= \frac{b \cdot \tfrac{a^2}{\abs{x_0}+\sqrt{x_0^2-a^2}}}{c}
\xrightarrow{x_0\to\infty} 0. \qedhere
\]
\end{proof}

\begin{notion}
Фокально-директрисное свойство~\eqref{21.1.focal} справедливо и для гиперболы с директрисами $x = \pm\dfrac{a^2}{c}$ и $\varepsilon > 1$.
\end{notion}

\section{Парабола}

\begin{defn}
\textit{Параболой} называется геометрическое место точек плоскости, равноудалённых от фиксированной точки $F$ (\textit{фокуса}) и фиксированной прямой $d$ (\textit{директрисы}):
\begin{equation}\label{21.3.def}
  \rho(P, F) = \rho(P, d).
\end{equation}
\end{defn}

Поместим фокус в точку $F\!\left(\dfrac{p}{2}, 0\right)$, $p > 0$, а директрису $d\colon x = -\dfrac{p}{2}$. Из условия~\eqref{21.3.def}:
\[
\sqrt{\!\left(x-\tfrac{p}{2}\right)^{\!2}+y^2} = x + \tfrac{p}{2},
\]
возводя в квадрат и упрощая, получим

\begin{equation}\label{21.3.canon}
  y^2 = 2px, \quad p > 0
\end{equation}
--- \textit{каноническое уравнение параболы}. Величина $p$ называется \textit{параметром параболы}.

\begin{notion}
Эксцентриситет параболы равен $\varepsilon = 1$, что согласуется с фокально-директрисным определением: отношение расстояния до фокуса к расстоянию до директрисы тождественно равно $1$.
\end{notion}

\section{Приведение общего уравнения к каноническому виду}

Общее уравнение~\eqref{21.0.general} приводится к каноническому виду поворотом и параллельным переносом. Запишем~\eqref{21.0.general} в матричном виде:

\begin{equation}\label{21.4.matrix}
  \vv{x}^{\top} \mathbf{A} \vv{x} + \vv{b}^{\top}\vv{x} + F = 0,
\end{equation}
где $\vv{x} = (x,y)^{\top}$,
\[
\mathbf{A} = \begin{pmatrix} A & B/2 \\ B/2 & C \end{pmatrix}, \quad \vv{b} = \begin{pmatrix} D \\ E \end{pmatrix}.
\]

\begin{enumerate}
  \item \textbf{Поворот.} Матрица $\mathbf{A}$ симметрична, поэтому ортогонально диагонализируема: $\mathbf{A} = \mathbf{Q}\,\mathrm{diag}(\lambda_1,\lambda_2)\,\mathbf{Q}^{\top}$. Замена $\vv{x} = \mathbf{Q}\vv{x}'$ (поворот на угол $\theta$, определяемый собственными векторами) обращает в нуль смешанный член $Bxy$.
  \item \textbf{Параллельный перенос.} После поворота уравнение принимает вид $\lambda_1 x'^2 + \lambda_2 y'^2 + D'x' + E'y' + F' = 0$. Выделяя полные квадраты по ненулевым $\lambda_k$, сдвигом начала координат убираем линейные члены.
\end{enumerate}

Тип кривой определяется знаками собственных значений $\lambda_1, \lambda_2$ матрицы $\mathbf{A}$ и знаком полного определителя
\begin{equation}
  \Delta = \begin{vmatrix} A & B/2 & D/2 \\ B/2 & C & E/2 \\ D/2 & E/2 & F \end{vmatrix}\!:
\end{equation}
\begin{itemize}
  \item $\lambda_1 \lambda_2 > 0$, $\Delta \neq 0$ --- \textbf{эллипс} (или окружность при $A=C$, $B=0$);
  \item $\lambda_1 \lambda_2 < 0$, $\Delta \neq 0$ --- \textbf{гипербола};
  \item $\lambda_1 \lambda_2 = 0$ (одно собственное значение равно нулю), $\Delta \neq 0$ --- \textbf{парабола};
  \item $\Delta = 0$ --- вырожденный случай (пара прямых, точка, пустое множество).
\end{itemize}

\section{Оптические свойства кривых второго порядка}

\begin{thm}[Отражательное свойство эллипса]
Касательная к эллипсу в любой точке $P$ образует равные углы с фокальными радиусами $PF_1$ и $PF_2$.
\end{thm}

\begin{proof}
Уравнение касательной к эллипсу~\eqref{21.1.canon} в точке $P(x_0, y_0)$:
\begin{equation}\label{21.5.tang_ell}
  \frac{x_0 x}{a^2} + \frac{y_0 y}{b^2} = 1.
\end{equation}
Нормаль к касательной в точке $P$ имеет направляющий вектор $\vv{n} = \left(\dfrac{x_0}{a^2},\, \dfrac{y_0}{b^2}\right)$. Прямое вычисление показывает, что косинусы углов между $\vv{n}$ и векторами $\vv{PF_1}$, $\vv{PF_2}$ совпадают, откуда следует равенство углов падения и отражения.
\end{proof}

\begin{thm}[Отражательное свойство параболы]
Касательная к параболе в любой точке $P$ образует равные углы с фокальным радиусом $PF$ и прямой, параллельной оси симметрии параболы.
\end{thm}

\begin{proof}
Уравнение касательной к параболе~\eqref{21.3.canon} в точке $P(x_0, y_0)$:
\begin{equation}\label{21.5.tang_par}
  y_0 y = p(x + x_0).
\end{equation}
Вектор нормали $\vv{n} = (-p,\, y_0)$. Угол между $\vv{n}$ и вектором $\vv{PF} = \left(\tfrac{p}{2}-x_0,\,-y_0\right)$ равен углу между $\vv{n}$ и вектором $(1, 0)$ (направление оси), что проверяется непосредственным вычислением косинусов.
\end{proof}